\section{Producing the selected finite-volume transfer data}\label{ex-centered-freight-producer}

The centered estimate uses a finite-volume representation at every period and volume. On the public construction route these data are produced on the same selected history: the run chooses one terminal voice, then reconstructs compatible voices at all earlier scales. This section isolates that producer and explains the two algebraic steps that turn the reconstruction into transfer data. The tariff and one-step interpolation estimates remain the quantitative inputs of the selected history; the proof below does not replace them by an unpriced compactness argument.

\subsection{One-step response and its finite-dimensional assembly}

Fix a volume, period, and adjacent scales $m,m+1$. Let $d$ be the chosen grade, let $\ell_m\geq1$ be the block side, let $F$ be its finite-dimensional complex coefficient space, and let $s\geq1$ be the number of paid columns at scale $m+1$. Write $\mu_F(u)=\sum_a|u_a|$ for its coefficient $\ell^1$ norm, where $|\cdot|$ is the ordinary complex modulus, and let $\operatorname{core}(u)$ erase the vacuum coordinate. Every frame coordinate has degree between $0$ and $d$. The scale dilation $\operatorname{vowel}_m$ multiplies coordinate $a$ by $\ell_m^{d-\deg(a)}$, so
\[
 \mu_F(u)\leq\mu_F(\operatorname{vowel}_m u)\leq\ell_m^d\mu_F(u).
\]
A voice at scale $m+1$ consists of a density family $f_v$, a nonnegative charge, and an index isomorphism $\iota:F\to F$. The selected voice certificate requires $\iota$ to commute with each motion on $F$. The family is complex-linear in $v$ and its motion covariance is semilinear according to the independent Boolean conjugation flag in $\mathrm{Motion}=\mathrm{Signed}\times\mathrm{Bool}$; it is not determined by coordinate orientation. For each $v$, the selected voice certificate also supplies the format and remainder bounds; the associated finite-volume integrability lemma gives $f_v\in L^1$ of the normalized prior. We do not use continuity of the densities. The selected transition defines the Wilson push
\[
  (\mathsf P_m f)(U)=\ell_m^d\int f(V)\,K_m(U,dV),
\]
where $K_m(U,dV)$ is the one-step Wilson transition density against normalized Haar measure and $\ell_m$ is the block side. Thus $\mathsf P_m$ includes the grade-scaling factor $\ell_m^d$; it is not a probability-preserving Markov operator. The selected transition density is jointly measurable, nonnegative and uniformly bounded in both configurations, and its integral in the coarse configuration equals one for every fine configuration. These properties give slice integrability: for every output $U$, the transition-weighted section of a prior-integrable input is integrable, and the resulting pushed function is measurable. This is the conditional-kernel conclusion attached to the lawful transition, not a consequence of finite total mass alone. The transition commutes with the selected signed motions. These properties make the push well-defined and linear, and for each signed motion $g$,
\[
  g\cdot(\mathsf P_m f_v)=\mathsf P_m f_{g\cdot v}.
\]
Indeed, both identities follow first for finite linear combinations from linearity of the integral; the motion identity follows by the change of variables preserving the Wilson transition density. These identities are exact, not limiting statements.

The selected ticket's one-column estimates have three exact forms. For each nonzero frame axis $e_a$, it supplies a coarse coefficient whose density is the pushed fine density, whose orbit-cost-weighted lead error is at most $\eta\,\|\operatorname{core}(e_a)\|_1$, whose activity residual price is at most $C\,\|\operatorname{vowel}_m(\operatorname{core}(e_a))\|_1$, and whose finite gate-crossing certificate costs $t$. For each regular supported term of mass $u$, both coupon prices are $\eta u$ and the gate-crossing cost is $2u$. For each forced term of mass $u$ on a marked wing site, both coupon prices are $\eta u/\kappa$ and the gate-crossing cost is $2u$, where $0<\kappa=\sqrt{\text{forcedAllowance}_{m+1}}\leq1$ uses the fine level. The output remainder is measured with the separate coarse-level forced allowance. Each such coupon includes a format bound, the exact density or activity identity, an orbit-cost lead bound, a finite response expansion, and its paid gate-crossing estimate through the selected door. These are the exact hypotheses used below; we do not infer their numerical bounds from compactness.

On the public route these coupons are produced together for the selected adjacent history. A held spring supplies an interval $[x,y]$, a parameter $z$ in that interval, and a local realization $K(z)$. The two selected fields are the fine representation $K(z)$ and the coarse result of that same spring at $z$. The gate is the composed local comparison between those fields; the same spring and parameter give the climb witness. No strict inequality between $x$ and $y$ is needed. For a grade $d$ with $d+3$ no larger than the coarse clock, the ticket's wing contains the principal-field tags indexed by supports containing the origin and the marked keeper tags. It is finite because the finite-volume lattice has finitely many supports and keeper tags.

Write $E_{m+1}$ for the support exponent at the fine level, $\vartheta_{m+1}$ for its profile title, and $b_m$ for the block side. Denote the coarse and fine clock values by $J_m^{\mathrm c}=J(m)$ and $J_m^{\mathrm f}=J(m+1)$. The raw coupon prices are $\eta_m=\operatorname{rate}(b_m)$ for the lead and $C_m=E(J_m^{\mathrm f})$ for the residual, where $E$ is the selected schedule. The exact total gate-crossing allowance is
\[
 t_m=2e^{E_{m+1}}\vartheta_{m+1}.
\]
The selected clock interpolation chooses one heading for the whole wing. Its three coupon families are exactly the axis, regular-activity, and forced-activity cases displayed above. In the raw ticket, the axis lead price is $\eta_m\mu_F(\operatorname{core}e_a)$, the residual price is $C_m\mu_F(\operatorname{vowel}_m(\operatorname{core}e_a))$, and the crossing bill is $t_m$. A regular support of mass $u$ has equal lead and residual price $\eta_m u$ and bill $2u$; a marked forced support of mass $u$ has equal price $\eta_m u/\kappa$ and bill $2u$. These are the three fields of the finite coupon certificate. The origin unit is supplied by its exact unit coupon.

The coupon-production step is conditional on its field estimates. Assume the selected spring supplies the two adjacent fields at the same parameter and a sound local section; assume the heading forecast and cap satisfy the two-sided map bound, the principal-field estimate has the displayed axis prices, each supported regular field has the displayed regular prices, and each slot-fitted forced field at a marked keeper site has the displayed forced prices. Then one heading and one ticket for the whole finite wing are obtained by applying those three estimates to their respective terms and collecting the resulting certificates. The heading forecast comes from the calibrated flight matrix, and the selected schedule-price bound supplies the cap's triangular-price estimate. The axis, regular, and forced estimates are separate hypotheses; no one of them follows from another. This is the precise conditional finite coupon package used here. It leaves the derivation of those analytic field estimates and the selected schedule-price data to their upstream producers.

More precisely, the source's finite interpolation lemma has this conditional interface: when the two fields at the selected spring parameter, slot table, motion/rate data, schedule-price bound, and keeper are those of the selected history, and $d+3\leq J_m^{\mathrm c}$, it returns one heading and a ticket for every wing containing the principal-field tags just described. The ticket has axis residual price $C_m$, lead price $\eta_m$, and crossing allowance $t_m$, with the regular and forced prices just listed. The heading and all prices are common to that whole wing. Substituting the finite wing formed by those tags together with the keeper's marked tags gives the ticket used by the public run. The interface and its three analytic coupon estimates are the precise source inputs here; this conditional assembly does not reprove their field-level derivation.

The ticket prices $C_m,\eta_m$ are then checked against the one-step reserves. Let $T_m:F\to F$ be the selected heading's coefficient isomorphism. The purchase inequalities are
\[
 \mu_F(T_m v)\leq C_m\mu_F(v),\qquad
 \mu_F(T_m^{-1}v)\leq C_m\mu_F(v).
\]
The same selected clock purchase pays the ticket through
\[
 \begin{aligned}
 1&\leq C_m, & 0&\leq\eta_m,\\
 \eta_m(1+s)\operatorname{orbitCost}(F)&\leq\operatorname{barunit}(b_m),\\
 (C_m+\eta_m s)\operatorname{orbitCost}(F)&\leq W_{m+1}.
 \end{aligned}
\]
where $s$ is the number of paid columns and $W_{m+1}$ is the available fine-level column reserve. The one-column theorem combines the raw coupon with these purchases: after all $s$ columns its residual and lead allowances are $c_m=C_m+\eta_m s$ and $\epsilon_m=\eta_m(1+s)$, and its crossing bill is $t_m+2s$. These inequalities are the actual interpolation and funding checks consumed by the one-column construction; the proof does not choose a new adjacent pair or reset the history.

The reason the ticket pays these three kinds of term is the local decomposition of the marked field. For a fine coefficient $P$, first apply the local rack map: it preserves the coefficient coordinate and exponential entry, while its keeper supplies the forced-profile and remainder bounds. The density is then
\[
 (\text{exponential entry})\,\cdot(\text{local regular part})
   +\sum_{j}(\text{forced activity at }j)\,\cdot(\text{dressed field at }j).
\]
The regular summands use the regular coupons. A forced summand with zero coefficient vanishes. For a nonzero forced coefficient, its support is marked and belongs to the selected finite wing; its absolute mass uses the forced coupon after division by $\kappa$. The origin term uses the principal axis coupon. The source's remnant and exponential envelopes dominate both countable activity series absolutely; Tonelli's theorem applied to the resulting absolute majorant proves integrability and permits summing the densities term by term. The coupon-summation identity adds the costs and errors, so the three displayed prices give the one-column response and preserve the exact pushed density.

The local selected-history ticket is a finite set of marked insertions with a linear cap. It is paid with nonnegative allowances $C,\eta,t$ and yields an endpoint response on every coefficient of an incoming voice. In the normalization used by the source, the response allowances after accounting for all $s$ columns are
\[
  c=C+\eta s,\qquad \epsilon=\eta(1+s),\qquad t+2s.
\]
For each frame coordinate $e_a$, the one-column estimate gives a marked polynomial $Z_a$ for $j_a=\iota^{-1}(e_a)$. Applying a selected motion $g$ gives $Z_{g,a}$ with density $g\cdot(\mathsf P_m f_{j_a})=\mathsf P_m f_{\iota^{-1}(g\cdot e_a)}$, lead error at most $\epsilon\,\mu_F(\operatorname{core}e_a)$, residual activity at most $cD_a$, and charged cost at most $t+2s$, where
\[
 D_a=\mu_F\!\left(\operatorname{vowel}_m(\operatorname{core}e_a)\right).
\]
Let $\mathcal G=\mathrm{Motion}=\mathrm{Signed}\times\mathrm{Bool}$ be the finite source motion set. Write $g=(\pi,\epsilon)$, where $\pi$ is a signed coordinate permutation and $\epsilon\in\{0,1\}$ is the independent conjugation flag. Let $\rho_g$ be its (possibly semilinear) action on $F$ and let $\sigma_g$ be its scalar action, identity for $\epsilon=0$ and complex conjugation for $\epsilon=1$. Then the actual averaging chips are
\[
 \chi(v;g,a)=\frac{1}{|\mathcal G|}\,
       \sigma_g\!\left((\rho_g^{-1}\iota v)_a\right),
 \qquad (g,a)\in\mathcal G\times\{1,\ldots,q\},
\]
where $q=\dim_{\mathbb C}F$. Applying the inverse motion, taking its $a$th coordinate, and then applying the matching scalar action makes $v\mapsto\chi(v;g,a)$ complex-linear, including when $\epsilon=1$. Averaging the reindexed coordinate expansions gives
\[
 \iota(v)=\sum_{g,a}\chi(v;g,a)\,\rho_g e_a,
 \qquad
 v=\sum_{g,a}\chi(v;g,a)\,\iota^{-1}(\rho_g e_a).
\]
For any selected $h$, reindexing $g=hq'$ gives
\[
 \chi(\rho_hv;hg,a)=\sigma_h(\chi(v;g,a)),
\]
so the assembled response has the required semilinear motion covariance. Thus conjugation comes from the independent flag, not from whether the signed coordinate permutation reverses orientation. The action preserves the grade and commutes with the core, but it need not permute the coordinate axes. Define the finite frame orbit price by
\[
 \operatorname{orbitCost}(F)
 :=1+\sum_{g\in\mathcal G}\sum_{a=1}^{q}\mu_F(\rho_g e_a).
\]
For any $x=\sum_a x_a e_a$, axis expansion and the triangle inequality give
\[
 \mu_F(\rho_g x)\leq\sum_a|x_a|\mu_F(\rho_g e_a)
 \leq\operatorname{orbitCost}(F)\mu_F(x).
\]
Apply this bound to $\rho_g^{-1}$ and to $x=\operatorname{core}(\iota v)$ or $x=\operatorname{vowel}_m(\operatorname{core}(\iota v))$. The scalar actions preserve modulus, and both core and vowel commute with the motion. Averaging over $g$ in the displayed chip formula gives the two estimates
\[
 \sum_{(g,a)}|\chi(v;g,a)|\,\mu_F(\operatorname{core}e_a)
 \leq \operatorname{orbitCost}(F)\mu_F(\operatorname{core}(\iota v)),
\]
\[
 \sum_{(g,a)}|\chi(v;g,a)|D_a
 \leq \operatorname{orbitCost}(F)\mu_F\!\left(\operatorname{vowel}_m(\operatorname{core}(\iota v))\right).
\]
They are the finite tally inequalities obtained by summing the absolute chip weights against the axis and dilated-axis costs. Therefore the weighted sums of lead errors and residual activities are at most $\epsilon\operatorname{orbitCost}(F)\mu_F(\operatorname{core}(\iota v))$ and $c\operatorname{orbitCost}(F)\mu_F(\operatorname{vowel}_m(\operatorname{core}(\iota v)))$, respectively; the corresponding finite tally of gate-crossing bills is the purchased wing allowance. These are the exact finite sums paid by the selected column and orbit-cost budgets. The estimates hold after each signed motion because the marked ticket is closed under those motions and the Wilson transition is motion covariant.
Here is the finite assembly step. Let $\iota:F\to F$ be the voice's coefficient-index isomorphism, and let $\mathcal J$ be the finite tally of pairs $(g,a)$ consisting of a selected signed motion $g$ and a coordinate index $a$. Put $e_a$ for the coordinate vector. Let $j_a=\iota^{-1}(e_a)$ and $v_{g,a}=g\cdot j_a=\iota^{-1}(g\cdot e_a)$; the equality uses equivariance of the index map. The frame interpolation coefficients $\chi(v;g,a)$ are the coefficients of the finite frame expansion of $\iota(v)$, so
\[
 \iota(v)=\sum_{(g,a)\in\mathcal J}\chi(v;g,a)(g\cdot e_a),\qquad
 v=\sum_{(g,a)\in\mathcal J}\chi(v;g,a)v_{g,a}.
\]
For each coordinate $a$, first take the one-column response at $j_a$, then apply the selected signed motion $g$ to it. The resulting moved response $Z_{g,a}$ has fine density $g\cdot f_{j_a}$ and, by covariance of the transition, coarse density $g\cdot(\mathsf P_m f_{j_a})=\mathsf P_m f_{v_{g,a}}$. Its scalar lead is one term in the reconstructed coordinate, and its forced profile is the corresponding moved profile. Assemble them by the finite sum
\[
 Z(v)=\sum_{(g,a)\in\mathcal J}\chi(v;g,a)Z_{g,a}.
\]
The assembled scalar lead and forced profile are the same finite weighted sums of their one-column counterparts. By the exact one-column density identity, linearity of the push, and the second frame identity,
\begin{align*}
 \operatorname{density}(Z(v))
 &=\sum_{(g,a)}\chi(v;g,a)g\cdot(\mathsf P_m f_{j_a})\\
 &=\mathsf P_m\Bigl(\sum_{(g,a)}\chi(v;g,a)f_{v_{g,a}}\Bigr)
 =\mathsf P_m f_v.
\end{align*}
All coefficient and residual bounds are summed with $|\chi(v;g,a)|$. The frame interpolation inequalities are proved by the finite frame identities: the lead error is bounded by $\epsilon\,\operatorname{orbitCost}(F)\|\operatorname{core}(\iota v)\|_1$, and the residual by $c\,\operatorname{orbitCost}(F)\|\operatorname{vowel}_m(\operatorname{core}(\iota v))\|_1$. The first is the triangle inequality followed by the frame orbit-cost bound; the second follows by summing the coordinate bounds and the vowel size inequality. The paid ticket bounds the total gate-crossing costs by the corresponding finite-wing budget. Thus no infinite series or exchange of an infinite sum with the transition integral occurs.
The one-step direction estimates give the cone bounds; here is the numerical argument. Put $L=b_m$, let $J_f$ and $J_c$ be the fine and coarse clocks with $J_f\leq J_c+1$, and set $C=W.K(J_f)$, $D_f=W.K(J_f+1)$, $D_c=W.K(J_c+1)$. Let $s_f$ be the fine column count, $s_c$ the coarse column count, and $e=\operatorname{barunit}(L)$. Write $u$ and $v$ for the coarse and fine frame coordinates, $r$ and $t$ for their charges, and $y=\operatorname{vowel}_L(v)$. The direction certificate provides a coefficient isomorphism $B$ such that
\[
 \|Bx\|_1,\ \|B^{-1}x\|_1\leq C\|x\|_1,
 \qquad B(\operatorname{core}x)=\operatorname{core}(Bx),
\]
\[
 \|u-By\|_1\leq e(\|\operatorname{core}v\|_1+t),
 \qquad r\leq C\|\operatorname{core}y\|_1+et,
 \qquad t\leq s_f\|\operatorname{core}v\|_1.
\]
The selected price growth gives $1\leq C\leq D_f$ and, by the slow-step condition, $C\leq D_c$. The bank certificate says
\[
 100000D_f^6(1+s_f)^2\leq \operatorname{puffs}(L),\qquad
 e=\operatorname{puffs}(L)^{-3}.
\]
In particular $\operatorname{puffs}(L)\geq1$ and $2C(1+s_f)\leq\operatorname{puffs}(L)$, so $2Ce(1+s_f)\leq\operatorname{puffs}(L)^{-2}\leq1$ and $e\geq0$. The fine room certificate gives $64D_f^2\leq s_f$, and the coarse room certificate gives $64D_c^2\leq s_c$. Hence $2C\leq s_f$ and $4C^2\leq s_c$. Since $\|\operatorname{core}y\|_1=\|\operatorname{vowel}_L(\operatorname{core}v)\|_1\geq\|\operatorname{core}v\|_1$, the incoming charge bound implies
\[
 2Ce(\|\operatorname{core}v\|_1+t)
 \leq 2Ce(1+s_f)\|\operatorname{core}v\|_1
 \leq \|\operatorname{core}y\|_1.
\]
Consequently $2C\|u-By\|_1\leq\|\operatorname{core}y\|_1$. Put $\delta=\|u-By\|_1$. Then $C\delta\leq\|y\|_1/2$, and the inverse norm bound gives
\[
 \|y\|_1=\|B^{-1}(By)\|_1
 \leq C\|By\|_1
 \leq C\|u\|_1+C\delta
 \leq C\|u\|_1+\tfrac12\|y\|_1,
\]
so $\|y\|_1\leq2C\|u\|_1$. In the other direction,
\[
 \|u\|_1\leq\|By\|_1+\delta
 \leq C\|y\|_1+\frac{1}{2C}\|y\|_1
 \leq2C\|y\|_1,
\]
since $C\geq1$. Apply the same absorption argument to the core vectors; the core projection contracts the $\ell^1$ norm and commutes with $B$. It yields
\[
 \|\operatorname{core}y\|_1\leq2C\|\operatorname{core}u\|_1.
\]
It follows, using $\|v\|_1\leq\|y\|_1\leq L^d\|v\|_1$, that
\[
 \|v\|_1\leq s_f\|u\|_1,\qquad
 \|u\|_1\leq s_f L^d\|v\|_1.
\]
Also $et\leq C\|\operatorname{core}y\|_1$, by the displayed bank estimate, so
\[
 r\leq2C\|\operatorname{core}y\|_1
   \leq4C^2\|\operatorname{core}u\|_1
   \leq s_c\|\operatorname{core}u\|_1.
\]
These are exactly the charge, fine-to-coarse, and coarse-to-fine bounds used below. In particular, the two norm inequalities are full coefficient norms; the core norm occurs only in the charge estimate.

The response defines a linear map $Q_m:F\to F$ by its coarse reference coordinate. The first cone inequality gives $\|\iota(v)\|_1\leq s_f\|Q_m v\|_1$. If $Q_m v=0$, then $\iota(v)=0$; since $\iota$ is an isomorphism, $v=0$. Thus $Q_m$ is injective and, since $F$ is finite dimensional, bijective. The one-step output stores this isomorphism as the new voice index. The density family is not precomposed by its inverse. Across several scales, the full-norm cone inequality is iterated from the reference to the terminal frame; the exact recurrence and terminal isomorphism give the single global transfer normalization used below.

\begin{lemma}[Finite backward reconstruction]\label{ex-freight-backward}
Let $t\leq N$. Suppose a selected finite-volume history has (i) a terminal voice which agrees with the calibrated top coefficient frame and satisfies the terminal reflection condition; (ii) a motion-covariant Wilson push between each adjacent pair of scales whose transition density has the boundedness, measurability and mass properties stated above; (iii) at every $m\in[t,N)$ a selected door, a sound local section, and a finite paid ticket whose one-step response satisfies the column estimates and cap conditions described above; and (iv) the room and bank bounds needed to pay the cumulative responses. Then there is a sequence of voices $V_m$, $t\leq m\leq N$, such that the terminal voice is the selected one, every adjacent density satisfies
\[
 \operatorname{density}(V_m(v))=\mathsf P_m\operatorname{density}(V_{m+1}(v)),
\]
and each scale has the selected reality, motion, and coefficient-cone bounds. The pushed coefficient maps define a linear transfer system, normalized by the inverse of its reference coefficient map.
\end{lemma}

\begin{proof}
Use downward induction. At $m=N$, take the supplied terminal voice. Suppose $V_{m+1}$ has been constructed and has the required integrability, reality, and motion properties. Apply the one-step response assembly above to its paid ticket. The density identity follows from the exact finite linear assembly. The local cap gives the response and residual bounds at scale $m$; the paid room and bank inequalities keep the accumulated expenditure within its assigned budget. The terminal reflection condition propagates because the Wilson push is invariant under reflection and the ticket is real on the selected real subspace. The coefficient-cone estimate propagates by the one-step paid inequality. The one-step construction packages its own invertible index map with its density family as a voice; we do not alter the family by precomposing with the inverse. This completes one induction step. There are only $N-t$ steps, so finite downward induction constructs every $V_m$ and all the claimed identities. Finally, if the reference readout of a coefficient $v$ is zero, the full-norm cone inequality propagates zero from scale $t$ to scale $N$. The terminal voice index is an isomorphism, so $v=0$. The reference map is injective and hence bijective on $F$; its inverse gives the one global normalization in the conclusion.\end{proof}

\subsection{The transfer map and its normalization}

For the unnormalized music just constructed, write $x_m(v)$ for the coefficient index of its voice at scale $m$. The full-norm cone estimates at each adjacent pair have the form
\[
 \|x_{m+1}(v)\|_1\leq s_{m+1}\|x_m(v)\|_1,
 \qquad
 \|x_m(v)\|_1\leq s_{m+1}\ell_m^d\|x_{m+1}(v)\|_1,
\]
where $s_{m+1}$ is the paid column count and $\ell_m$ the block side. Let $Q_t(v)=x_t(v)$ be the reference readout. If $Q_t(v)=0$, the first inequality propagates zero successively through every scale up to $N$. At the terminal scale $x_N$ is an isomorphism, so $v=0$. Thus $Q_t$ is injective, and finite dimensionality makes it bijective. This is the sole global normalization. The transfer data record the unnormalized music entries and set their terminal coordinate map to the reference readout $Q_t$. When viewed on external labels $w$, every scale uses the same reparametrization $V_m(Q_t^{-1}w)$; no scale gets an independently fitted inverse.
The selected public supplier applies this construction for every grade and every pair $(N,p)$ beyond the grade reference. Its first finitary step is the selected endpoint: the terminal frame is obtained from the calibrated section and the keeper chosen by the nested history. The endpoint is on top because the selected Wilson trajectory represents the terminal field and the paid terminal receipt gives the exact coefficient identities. Its reflection condition follows from the real section and the terminal keeper. The selected climb provides the door and the finite ticket at each adjacent pair. The ticket's allowance is the scale's tariff, and the paid room and bank bounds control the total cost. Applying Lemma~\ref{ex-freight-backward} and then the reference normalization gives the transfer data required by the centered-volume calculation.

The resulting ordinary mathematical object is a music of compatible voices with a selected terminal frame, the adjacent finite-wing estimates, and a normalized transfer map. In the center argument, the separate terminal birth polynomial has coefficients given by finite products of the named observables. It is continuous on compact configuration space and hence integrable. The voice and its transition identities supply the individual insertion profiles used in the served representation. The full birth polynomial is propagated separately; its higher-mode coefficients belong to the census construction, not to a single voice. The backward integral identity then relates each propagated mode to the corresponding terminal moment. All later bounds concern normalized finite-volume expectations.

\subsection{Scope of the producer statement}

The preceding induction proves the finite reconstruction from its exact selected-history ticket and terminal inputs. In the public run, the schedule is chosen from the purchased trajectory so that the grade reference lies before every requested terminal level. The selected climb and terminal receipt are retained on this same trajectory; they are not assembled by selecting a new history for each period. The consumer then uses the resulting transfer system to form the terminal birth and the root estimate.

The quantitative production of the ticket's nonlinear coefficient estimates is a distinct finite-volume step. The selected public construction supplies it through the field interpolation, mode-slot, and shipment records; this chapter only states the numerical inequalities that the one-step theorem consumes. Likewise, uniformity over every coefficient in a bounded grade ball must be established from those records before the conditional centered theorem becomes unconditional. The evidence file records the exact source chain and distinguishes the finite backward assembly proved here from these upstream quantitative estimates.
